\documentclass[12pt]{article}
% Page setup
\usepackage[margin=1in]{geometry}
\usepackage{parskip}
\usepackage{titling}
\setlength{\droptitle}{-6em} 

% Math packages
\usepackage{amsmath, amsthm, amssymb, amsfonts}
\usepackage{mathtools}
\makeatletter
\g@addto@macro\th@plain{\normalfont} % remove italics from theorems (plain style)
\makeatother

% Commutative diagrams and proofs
\usepackage{tikz-cd}
\usepackage{proof}
\usepackage{stmaryrd}

% Unordered lists
\usepackage{enumitem}
\newlist{related}{itemize}{1}
\setlist[related]{label=\textbullet}
\newlist{distinct}{itemize}{1}
\setlist[distinct]{label=--}


% Theorem environments
\newtheorem{theorem}{Theorem}[section]
\newtheorem{lemma}[theorem]{Lemma}
\newtheorem{proposition}[theorem]{Proposition}
\newtheorem{corollary}[theorem]{Corollary}
\newtheorem{fact}[theorem]{Fact}
\newtheorem{remark}[theorem]{Remark}

\theoremstyle{definition}
\newtheorem{definition}[theorem]{Definition}
\newtheorem{example}[theorem]{Example}

% Useful commands for category theory
\newcommand{\alp}{\mathbf{Alph}}
\newcommand{\mc}{\mathcal}
\newcommand{\ini}{\operatorname{ini}}
\newcommand{\ter}{\operatorname{ter}}
\newcommand{\cod}{\operatorname{cod}}
\newcommand{\dom}{\operatorname{dom}}

\newcommand{\List}{\operatorname{List}}
\newcommand{\lop}{\operatorname{[ \ ]}}
\newcommand{\len}{\operatorname{len}}
\newcommand{\type}{\operatorname{type}}

\newcommand{\Exc}{\operatorname{Exception}}

\newcommand{\Kl}{\operatorname{Kl}}
\newcommand{\plus}{\operatorname{+}}

\newcommand{\Cat}{\mathcal{C}at}
\newcommand{\Vect}{\mathrm{Vect}_\mathbf{k}}
\newcommand{\Hom}{\operatorname{Hom}}
\newcommand{\Nat}{\operatorname{Nat}}
\newcommand{\id}{\mathrm{id}}
\newcommand{\op}{\mathrm{op}}
\newcommand{\obj}{\operatorname{obj}}
\newcommand{\arr}{\operatorname{arr}}
\newcommand{\adj}{\operatorname{adj}}
\newcommand{\Par}{\operatorname{par}}

\newcommand{\coadj}{\operatorname{coadj}}
\newcommand{\colim}{\operatorname*{colim}}

% Arrows
\newcommand{\ra}{\rightarrow}
\newcommand{\from}{\leftarrow}
\newcommand{\To}{\Rightarrow}
\newcommand{\xto}{\xrightarrow}
\newcommand{\xfrom}{\xleftarrow}

\renewcommand{\labelitemi}{--} % Level 1: Dashes
\renewcommand{\labelitemii}{$\circ$} % Level 2: Circles
\renewcommand{\labelitemiii}{$\bullet$} % Level 3: Bullets (default)


\title{Category Theory out of Context}
\author{Alyson Mei}
\date{\today}


\begin{document}

\maketitle

%\tableofcontents


\section{Basics}

\paragraph{Base definitions}

\begin{definition}[Category] A \emph{category $\mc{C}$} consists of the following data:
\begin{enumerate}
    \item A collection of \emph{objects} $\obj(\mc{C})$;
    \item A collection of \emph{arrows} $\arr(\mc{C})$ of form $f: c \to d$ for $c, d: \obj(\mc{C})$ satisfying the \emph{composition}, \emph{identity} and \emph{associativity} axioms:
    % https://q.uiver.app/#q=WzAsMTEsWzEsMCwiZCJdLFsxLDEsImQiXSxbMCwwLCJjIl0sWzIsMCwiYyJdLFsyLDEsImMiXSxbMywxLCJkIl0sWzMsMCwiZCJdLFs0LDAsImMiXSxbNSwwLCJkIl0sWzUsMSwiZSJdLFs2LDEsInciXSxbMiwwLCJmIl0sWzAsMSwiZyJdLFsyLDEsImcgXFxjaXJjIGYiLDJdLFszLDQsIlxcaWRfYyIsMl0sWzQsNSwiZiIsMl0sWzMsNSwiZiIsMV0sWzMsNiwiZiJdLFs2LDUsIlxcaWRfZCJdLFs3LDgsImYiXSxbOCw5LCJnIiwxXSxbOSwxMCwiaCIsMl0sWzcsOSwiZyBcXGNpcmMgZiIsMl0sWzgsMTAsImggXFxjaXJjIGciXV0=
    \[\begin{tikzcd}
        c & d & c & d & c & d & \\
        & d & c & d && e & w
        \arrow["f", from=1-1, to=1-2]
        \arrow["{g \circ f}"', from=1-1, to=2-2]
        \arrow["g", from=1-2, to=2-2]
        \arrow["f", from=1-3, to=1-4]
        \arrow["{\id_c}"', from=1-3, to=2-3]
        \arrow["f"{description}, from=1-3, to=2-4]
        \arrow["{\id_d}", from=1-4, to=2-4]
        \arrow["f", from=1-5, to=1-6]
        \arrow["{g \circ f}"', from=1-5, to=2-6]
        \arrow["g"{description}, from=1-6, to=2-6]
        \arrow["{h \circ g}", from=1-6, to=2-7]
        \arrow["f"', from=2-3, to=2-4]
        \arrow["h"', from=2-6, to=2-7]
    \end{tikzcd}\text{.}\] 
\end{enumerate}
\end{definition}

\begin{remark}[Arrow-first approach]
    We identify objects with their identity arrows, so we may to assume that categories consist of arrows only. We use this viewpoint when convenient. 
\end{remark}

\begin{definition}[Functor]
    A functor $F: \mc{C} \to \mc{D}$ is a structure-preserving map between categories:
    \[
        F(c \xto{\id_c} c) = Fc \xto{\id_{Fc}} Fc, \qquad F(c \xto{f} d \xto{g} e) = Fc \xto{Ff} Fd \xto{Fg} Fe.
    \]
\end{definition}

\begin{definition}[Natural transformation] A \emph{natural transformation} $\alpha: F \to G$ between $F, G: \mc{C} \to \mc{D}$ is given by a collection of arrows 
\[
    (\alpha_c: Fc \to Gc \mid c: \obj(\mc{C})),
\]
that for each $c \xto{f} c': \arr(\mc{C})$ satisfies the 
naturality condition
% https://q.uiver.app/#q=WzAsNCxbMCwwLCJGYyJdLFsxLDAsIkdjIl0sWzEsMSwiR2MnIl0sWzAsMSwiRmMnIl0sWzAsMywiRmYiLDJdLFsxLDIsIkdmIl0sWzAsMSwiXFxhbHBoYV9jIl0sWzMsMiwiXFxhbHBoYV97Yyd9IiwyXV0=
\[\begin{tikzcd}
	Fc & Gc \\
	{Fc'} & {Gc'}
	\arrow["{\alpha_c}", from=1-1, to=1-2]
	\arrow["Ff"', from=1-1, to=2-1]
	\arrow["Gf", from=1-2, to=2-2]
	\arrow["{\alpha_{c'}}"', from=2-1, to=2-2]
\end{tikzcd}.\]
\end{definition}

\paragraph{Special objects and arrows}

\begin{definition}[Universal objects]
\hfill
\begin{distinct}
    \item \emph{Initial object} $\ini(\mc{C})$
    \item \emph{Terminal object} $\ter(\mc{C})$
    \item \emph{Universal object} is either initial or terminal.
\end{distinct}
\end{definition}

\begin{definition}[Monomorphisms and epimorphisms]
    
\end{definition}

\end{document}
